% Página de agradecimientos del trabajo de grado.

\section*{Agradecimientos}

En primer lugar, expresamos nuestra m\'as profunda gratitud a \textbf{Dios}, por haber sido la luz que ilumin\'o cada etapa de este proceso, por brindarnos la salud, la sabidur\'ia y la perseverancia necesarias para afrontar los retos que conlleva la culminaci\'on de un proyecto de esta magnitud. A \'El debemos las oportunidades que se nos presentaron a lo largo de la carrera y la fortaleza para no rendirnos en los momentos en los que el camino se torn\'o particularmente exigente.

Queremos expresar un agradecimiento especial a nuestro asesor, el ingeniero \textbf{Juan Pablo Ospina}, cuya gu\'ia metodol\'ogica, rigurosidad acad\'emica y disposici\'on permanente fueron determinantes para dar forma, estructura y profundidad a este trabajo. Sus observaciones precisas, su capacidad para se\~nalar oportunidades de mejora y su acompa\~namiento durante cada fase del proyecto se convirtieron en un referente invaluable para nuestro proceso de formaci\'on como profesionales.

De igual manera, extendemos nuestro reconocimiento al ingeniero \textbf{Camilo Rodr\'iguez}, quien en su rol de coasesor aport\'o una visi\'on cr\'itica y constructiva sobre los aspectos t\'ecnicos del proyecto, ayud\'andonos a tomar mejores decisiones de arquitectura, dise\~no y desarrollo. Su experiencia y su disponibilidad para escuchar nuestras inquietudes nos permitieron consolidar una propuesta s\'olida y coherente con las exigencias del entorno profesional actual.

Agradecemos tambi\'en al ingeniero \textbf{Camilo Salazar}, cuya orientaci\'on a lo largo de distintas asignaturas y espacios acad\'emicos contribuy\'o significativamente a fortalecer las bases conceptuales y pr\'acticas sobre las cuales se sustenta este proyecto. Su pasi\'on por la ense\~nanza y su compromiso con la formaci\'on de profesionales \'integros constituyen un ejemplo que llevaremos con nosotros m\'as all\'a de las aulas.

As\'i mismo, queremos reconocer al ingeniero \textbf{William Frasser}, cuyas ense\~nanzas, recomendaciones y aportes durante el desarrollo de la carrera resultaron fundamentales para comprender la importancia de la calidad, la planeaci\'on y la mejora continua en cualquier proyecto de ingenier\'ia. Su acompa\~namiento nos permiti\'o ampliar nuestra perspectiva y abordar los retos del proyecto con mayor madurez t\'ecnica y profesional.

Un agradecimiento profundo y especial merecen nuestras \textbf{familias}, pilar fundamental e insustituible a lo largo de toda nuestra formaci\'on. A nuestros padres, hermanos y dem\'as seres queridos, quienes nos brindaron su apoyo incondicional, su comprensi\'on durante las largas jornadas de estudio, su paciencia ante las ausencias y su aliento constante en los momentos de mayor exigencia, les debemos buena parte de lo que hoy somos. Sus sacrificios silenciosos, sus palabras de \'animo y la confianza depositada en cada uno de nosotros fueron, sin duda alguna, el sost\'en emocional que hizo posible llegar hasta este punto. Este logro no es solo nuestro: tambi\'en es de ellos, porque cada paso dado en esta carrera lleva impreso el cari\~no, el esfuerzo y la entrega de quienes nos acompa\~naron desde el hogar.

Por \'ultimo, pero no por ello menos importante, extendemos nuestro m\'as sincero agradecimiento a la \textbf{Universidad Sergio Arboleda}, instituci\'on que nos abri\'o sus puertas y nos brind\'o no solo una formaci\'on acad\'emica de alto nivel, sino tambi\'en los espacios, recursos, herramientas y oportunidades que hicieron posible la realizaci\'on de este trabajo de grado. A su cuerpo docente, administrativo y a todas las personas que de una u otra manera hicieron parte de nuestro proceso formativo, les expresamos nuestro reconocimiento y nuestra gratitud por haber sido parte fundamental de esta etapa que hoy concluye y que marca el inicio de nuevos retos profesionales.

%  Se debe dejar la nueva página para guardar el orden de las seciones
\newpage
